\chapter{Discussion}
\label{chap:discussion}

\section{What closes the theory--practice gap}

\begin{itemize}
  \item Interpret the principal result as a complete chain of compatible decisions, not solely as a new transition recurrence.
  \item Identify which workflow stages changed downstream results most strongly.
  \item Explain where explicit units, time intervals, projections, state metadata, and raster alignment prevented otherwise plausible but invalid analyses.
  \item Distinguish algorithm availability from practical usability through MoveApps interfaces, diagnostics, and artifacts.
\end{itemize}

\section{Scientific interpretation}

\begin{itemize}
  \item More environmentally plausible paths do not identify the animal's true missing route.
  \item Improved hidden-point probability indicates better predictive uncertainty under the evaluated model, not proof of a behavioral mechanism.
  \item State-specific kernels reflect the segmentation method and feature set as well as biological heterogeneity.
  \item Learned terrain weights are relative selection factors conditional on the kernel-defined availability and observed study design.
  \item Finite-interval conditioned UDs must not be presented automatically as long-term home range.
\end{itemize}

\section{Methodological trade-offs}

\begin{itemize}
  \item Flexible transition kernels can represent persistence, behavior, terrain, and barriers beyond standard Brownian bridges.
  \item Flexibility adds data requirements, identifiability questions, discretization choices, and computational cost.
  \item Criteria segmentation offers interpretability but depends on thresholds; HMM/BCPA methods discover structure but remain model- and scale-dependent.
  \item Detailed WorldCover improves local context but introduces remote-service and reference-year dependencies; binary long-range land/water annotation trades detail for robust coverage.
  \item Simple baselines remain preferable when data cannot support state-specific or terrain-specific parameters.
\end{itemize}

\section{Limitations}

\begin{itemize}
  \item Tracking error and non-random missingness are not fully modeled.
  \item Environmental layers differ from the animal observations in resolution, accuracy, and reference time.
  \item One UTM zone is unsuitable for some wide-ranging tracks; the fallback environmental mode is intentionally coarse.
  \item Rare states and rare terrain crossings yield unstable or neutral estimates.
  \item Grid, time, support, direction-bin, and boundary choices affect reachability and probability surfaces.
  \item Terrain-weight estimation currently precedes full automatic installation in every native mapping.
  \item Current result folders include exploratory runs from different configurations and must be frozen before final inference.
  \item GPU conclusions depend on hardware, memory layout, transfer costs, and the exact workload.
\end{itemize}

\section{Future work}

\begin{itemize}
  \item Integrate terrain-pair estimation and installation into one validated workflow stage.
  \item Quantify uncertainty in environmental annotation, behavioral states, kernels, and learned terrain weights jointly.
  \item Support robust long-range projections or adaptive spatial representations.
  \item Add hierarchical pooling across animals and principled treatment of rare states.
  \item Extend static environmental context with weather, currents, seasonality, and time-varying availability.
  \item Improve animated Leaflet explanations of endpoint conditioning and uncertainty without overwhelming users.
  \item Publish versioned experiment manifests linking every thesis figure and table to code, configuration, and raw outputs.
\end{itemize}
